\begin{courseunit}{II.1}
\chapter{Relative, interrogative, and indefinite pronouns}

\begin{learningobjectives}
  \item decline \Latin{qui, quæ, quod} and \Latin{quis, quid} without a chart;
  \item choose a relative pronoun's gender and number from its antecedent but its case from
        its own clause;
  \item distinguish a question, an ordinary relative clause, and a headless relative clause;
  \item recognize the indefinites most frequent in Scripture and the Missal.
\end{learningobjectives}

\section{Exact forms and rules}

A \Term{relative pronoun} typically introduces a clause that describes or identifies an
\Term{antecedent}, the noun or noun-equivalent to which it refers. An
\Term{interrogative pronoun} asks for missing information about a person or
thing; an \Term{indefinite pronoun}
refers without identifying a particular person or thing. A pronoun is
\Term{substantival} when it occupies a noun's slot (``who?''), but
\Term{adjectival} when it modifies a noun (``which offering?'').

\begin{grammarbox}{The relative pronoun: \Latin{qui, quæ, quod}}
\begin{center}
\begin{tabular}{@{}l*{3}{c}@{\qquad}*{3}{c}@{}}
\toprule
& \multicolumn{3}{c}{Singular} & \multicolumn{3}{c}{Plural}\\
& M & F & N & M & F & N\\ \midrule
Nom. & qui & quæ & quod & qui & quæ & quæ\\
Gen. & cuius & cuius & cuius & quorum & quarum & quorum\\
Dat. & cui & cui & cui & quibus & quibus & quibus\\
Acc. & quem & quam & quod & quos & quas & quæ\\
Abl. & quo & qua & quo & quibus & quibus & quibus\\
\bottomrule
\end{tabular}
\end{center}

The antecedent determines gender and number; the relative pronoun's own job in
its clause determines case. Thus in
\Latin{Hostiam quam offerimus}, \Latin{quam} is feminine singular because of
\Latin{hostiam}, and accusative because it is the object of \Latin{offerimus}.  A relative can
have no expressed antecedent. This is a \Term{headless relative}: the relative
clause itself includes the sense of the missing antecedent, so \Latin{quod
accepimus} means ``what [the thing which] we have received.''
\end{grammarbox}

\begin{grammarbox}{Questions and indefinites}
Substantival \Latin{quis? quid?} asks ``who? what?''; adjectival \Latin{qui? quæ? quod?}
asks ``which? what kind of?''  Except for nominative singular masculine \Latin{quis} and
nominative/accusative singular neuter \Latin{quid}, the interrogative declines like the
relative.  \Latin{uter, utra, utrum} asks ``which of two?'' and declines like \Latin{pulcher}.

\begin{tabularx}{\linewidth}{@{}P{1.35in}Y Y@{}}
\toprule
\textbf{Form} & \textbf{Force} & \textbf{Rule or warning}\\ \midrule
\Latin{aliquis, aliquid} & someone, something & after \Latin{si, nisi, ne, num}, the prefix
  normally drops: \Latin{si quis}, ``if anyone''\\
\Latin{quidam, quædam, quoddam} & a certain & the \Latin{-m} assimilates before \Latin{d}:
  \Latin{quendam, quandam}\\
\Latin{quisque, quidque} & each one, each thing & enclitic \Latin{-que} belongs to the pronoun\\
\Latin{quicumque, quæcumque, quodcumque} & whoever, whatever & decline only the first part\\
\Latin{nemo} & no one & nom. \Latin{nemo}, dat. \Latin{nemini}, acc.
  \Latin{neminem}; gen./abl. commonly \Latin{nullius/nullo}\\
\Latin{nihil} & nothing & indeclinable neuter in nom./acc.; oblique senses
  normally use \Latin{nullius rei, nulli rei, nulla re}; limited idioms use
  \Latin{nihili/nihilo}\\
\Latin{omnis; uterque; alius} & every; each of two; another & adjectives often used
  substantivally; \Latin{alius} has genitive \Latin{alius}, dative \Latin{alii}\\
\bottomrule
\end{tabularx}

An \Term{enclitic} is a short element attached to the end of the preceding
word rather than standing independently. Thus the \Latin{-que} of \Latin{quisque} is part of the
pronoun; it is not the coordinating conjunction \Latin{-que} attached later.
\end{grammarbox}

\section{Reading the forms in context}

\begin{workedexample}{an implied antecedent in the Gloria}{W.II.1}
\Latin{Qui sedes ad dexteram Patris, miserere nobis.}

\textbf{Analysis.}  \Latin{Qui} is nominative masculine singular.  Its antecedent is the
Christ addressed as an implied \Latin{tu}; it is the subject of second-person
\Latin{sedes}.  The imperative \Latin{miserere} governs dative \Latin{nobis}.

\textbf{Translation.}  ``You who sit at the Father's right hand, have mercy on us.''
\end{workedexample}

\begin{workedexample}{a headless relative in a Postcommunion}{W.II.2}
\Latin{Quod ore sumpsimus, Domine, pura mente capiamus.}

\textbf{Analysis.}  \Latin{Quod} is accusative neuter singular, the object within
\Latin{sumpsimus}; the entire clause \Latin{quod ore sumpsimus} is the object of hortatory
\Latin{capiamus}. A \Term{hortatory subjunctive} urges ``let us'' action.
The form \Latin{ore} is means and \Latin{pura mente} manner.

\textbf{Translation.}  ``Lord, may we take in with a pure mind what we have received by
mouth.''
\end{workedexample}

\begin{workedexample}{question and relative together (Psalm 115:12)}{W.II.3}
\Latin{Quid retribuam Domino pro omnibus quæ retribuit mihi?}

\textbf{Analysis.}  Interrogative \Latin{quid} is the neuter singular object of
\Latin{retribuam}.  Substantival \Latin{omnibus} means ``all the things''; relative
\Latin{quæ} is neuter plural because of that antecedent and accusative as the object of
\Latin{retribuit}.

\textbf{Translation.}  ``What shall I return to the Lord for all that he has given me?''
\end{workedexample}

\begin{workedexample}{a negative indefinite (John 14:6)}{W.II.4}
\Latin{Nemo venit ad Patrem, nisi per me.}

\textbf{Analysis.}  \Latin{Nemo} is nominative singular and the subject of \Latin{venit};
\Latin{nisi} excepts one route from the universal negative.

\textbf{Translation.}  ``No one comes to the Father except through me.''
\end{workedexample}

\section{Writing laboratory}

Copy the full relative paradigm twice.  On the first copy underline the three forms in
\Latin{-ius} or \Latin{-i}; on the second, circle every neuter nominative/accusative identity.
Then study, cover, and reproduce these printed phrases from memory: \Latin{qui tollis
peccata mundi}; \Latin{quod ore sumpsimus}; \Latin{pro omnibus quæ retribuit mihi}.
Label antecedent, case-giver, and clause function.

For controlled composition, combine each pair with a relative: ``We praise the Lamb.  The
Lamb takes away sins'' $\rightarrow$ \Latin{Laudamus Agnum, qui peccata tollit}; ``Receive the
offering.  We offer it'' $\rightarrow$ \Latin{Suscipe oblationem quam offerimus}.  Finally
write one question with \Latin{quid} and one condition with \Latin{si quis}.



\end{courseunit}

\begin{courseunit}{II.2}
\chapter{Comparison, adverbs, and numerals}

\begin{learningobjectives}
  \item form and decline regular and irregular comparatives and superlatives;
  \item identify both constructions of comparison and the ablative of degree;
  \item derive and compare adverbs;
  \item read the cardinal, ordinal, and distributive numerals needed in rubrics and lessons.
\end{learningobjectives}

\section{Comparison exactly stated}

The \Term{positive degree} names a quality without comparison (``holy''). The
\Term{comparative degree} presents a higher or lower degree (``holier''). The
\Term{superlative degree} presents the highest degree or, in some contexts, a
very high degree (``holiest'' or ``very holy''). An \Term{adverb} modifies a
verb, adjective, or another adverb rather than agreeing with a noun.

\begin{grammarbox}{Adjectives and adverbs}
A regular comparative adds \Latin{-ior} (masculine/feminine) and \Latin{-ius} (neuter) to the
adjective stem and declines as a third-declension consonant stem: genitive \Latin{-ioris},
dative \Latin{-iori}, accusative masculine/feminine \Latin{-iorem}, ablative singular
\Latin{-iore}, nominative/accusative neuter plural \Latin{-iora}.  A regular superlative adds
\Latin{-issimus, -a, -um}; stems in \Latin{-r} normally add \Latin{-rimus}
(\Latin{pulcher, pulcherrimus}); six adjectives in \Latin{-ilis} commonly add \Latin{-limus}
(including \Latin{facillimus}, \Latin{humillimus}, \Latin{simillimus}).

\begin{center}
\begin{tabular}{@{}llll@{}}
\toprule
Positive & Comparative & Superlative & Adverbs\\ \midrule
\Latin{bonus} & \Latin{melior, melius} & \Latin{optimus} &
  \Latin{bene, melius, optime}\\
\Latin{malus} & \Latin{peior, peius} & \Latin{pessimus} &
  \Latin{male, peius, pessime}\\
\Latin{magnus} & \Latin{maior, maius} & \Latin{maximus} &
  \Latin{magnopere, magis, maxime}\\
\Latin{parvus} & \Latin{minor, minus} & \Latin{minimus} &
  \Latin{paulum, minus, minime}\\
\Latin{multus} & \Latin{plus, pluris} n. sg.; \Latin{plures, plura} pl. & \Latin{plurimus} &
  \Latin{multum, plus, plurimum}\\
\bottomrule
\end{tabular}
\end{center}

\begin{center}
\begin{tabular}{@{}l*{2}{c}@{\qquad}*{2}{c}@{}}
\toprule
& \multicolumn{2}{c}{Singular} & \multicolumn{2}{c}{Plural}\\
& M/F & N & M/F & N\\ \midrule
Nom. & \Latin{-ior} & \Latin{-ius} & \Latin{-iores} & \Latin{-iora}\\
Gen. & \multicolumn{2}{c}{\Latin{-ioris}} & \multicolumn{2}{c}{\Latin{-iorum}}\\
Dat. & \multicolumn{2}{c}{\Latin{-iori}} & \multicolumn{2}{c}{\Latin{-ioribus}}\\
Acc. & \Latin{-iorem} & \Latin{-ius} & \Latin{-iores} & \Latin{-iora}\\
Abl. & \multicolumn{2}{c}{\Latin{-iore}} & \multicolumn{2}{c}{\Latin{-ioribus}}\\
\bottomrule
\end{tabular}
\end{center}

First/second-declension adjectives normally make adverbs in \Latin{-e}; third-declension
adjectives use \Latin{-iter} (or \Latin{-er} after \Latin{-ns}): \Latin{digne},
\Latin{fortiter}, \Latin{prudenter}. The comparative adverb is normally the
neuter singular comparative. The superlative adverb is formed from the
superlative-adjective stem plus \Latin{-e}: \Latin{sanctissime,
pulcherrime, facillime, optime}.
\end{grammarbox}

\begin{grammarbox}{Comparison and counting}
The \Term{standard of comparison} is the person, thing, or group against which
the comparative is measured. Parallel terms in the same clause normally use
\Latin{quam} plus the same case; \Latin{quam} may instead introduce an
elliptical clause whose noun takes its own function and case, as \Latin{quam
hæc dilectio [est]}. An ablative without \Latin{quam} may express the standard
when the compared word is nominative or accusative.
\Latin{Maior quam frater} and \Latin{maior fratre} both mean ``greater than the brother.''
An ablative such as \Latin{multo} states degree of difference: ``much greater.''  A superlative
may take a \Term{partitive genitive}, which names the whole group from which the
highest member is selected, or \Latin{ex/de} plus ablative.

\begin{tabularx}{\linewidth}{@{}lY@{\qquad}lY@{}}
\toprule
\textbf{Cardinal} & \textbf{Form} & \textbf{Ordinal} & \textbf{Form}\\ \midrule
1 & \Latin{unus, -a, -um} & 1st & \Latin{primus}\\
2 & \Latin{duo, duæ, duo} & 2nd & \Latin{secundus}\\
3 & \Latin{tres, tria} & 3rd & \Latin{tertius}\\
4--10 & \Latin{quattuor, quinque, sex, septem, octo, novem, decem} & 4th--10th &
\Latin{quartus, quintus, sextus, septimus, octavus, nonus, decimus}\\
100 & \Latin{centum} & 100th & \Latin{centesimus}\\
1000 & \Latin{mille}; plural \Latin{milia} + genitive & 1000th & \Latin{millesimus}\\
\bottomrule
\end{tabularx}

\Term{Cardinal numerals} answer ``how many''; \Term{ordinal numerals} mark
position in a sequence (first, second); \Term{distributive numerals} answer
``how many apiece''; and \Term{multiplicative adverbs} answer ``how many
times.'' \Latin{unus} has genitive \Latin{unius}, dative \Latin{uni}; \Latin{duo} and \Latin{tres}
decline.  Ordinals decline like \Latin{bonus}.  Distributives answer ``how many each'':
\Latin{singuli, bini, terni}.  Multiplicative adverbs include \Latin{semel, bis, ter}.

\begin{center}
\begin{tabular}{@{}llll@{}}
\toprule
& \Latin{unus, una, unum} & \Latin{duo, duæ, duo} & \Latin{tres, tres, tria}\\ \midrule
Nom. & \Latin{unus, una, unum} & \Latin{duo, duæ, duo} & \Latin{tres, tres, tria}\\
Gen. & \Latin{unius} & \Latin{duorum, duarum, duorum} & \Latin{trium}\\
Dat. & \Latin{uni} & \Latin{duobus, duabus, duobus} & \Latin{tribus}\\
Acc. & \Latin{unum, unam, unum} & \Latin{duos, duas, duo} & \Latin{tres, tres, tria}\\
Abl. & \Latin{uno, una, uno} & \Latin{duobus, duabus, duobus} & \Latin{tribus}\\
\bottomrule
\end{tabular}
\end{center}

Singular \Latin{mille} is indeclinable and directly modifies a noun.  Plural \Latin{milia}
is a neuter plural noun: nominative/accusative \Latin{milia}, genitive \Latin{milium},
dative/ablative \Latin{milibus}; the counted noun is genitive plural.
\end{grammarbox}

\section{Reading the forms in context}

\begin{workedexample}{a substantive superlative in the Gloria}{W.II.5}
\Latin{Tu solus Altissimus, Iesu Christe.}

\textbf{Analysis.}  \Latin{Altissimus} is nominative masculine singular superlative of
\Latin{altus} and is used substantivally as a title; \Latin{solus} limits \Latin{tu}.

\textbf{Translation.}  ``You alone are the Most High, Jesus Christ.''
\end{workedexample}

\begin{workedexample}{ablative standard (John 15:13)}{W.II.6}
\Latin{Maiorem hac dilectionem nemo habet.}

\textbf{Analysis.}  \Latin{maiorem} modifies accusative \Latin{dilectionem}.  Ablative
\Latin{hac} is a substantival demonstrative supplying the standard without \Latin{quam}, with
a second \Latin{dilectione} understood: ``greater than this [love].''  The printed
\Latin{dilectionem} is accusative and belongs with \Latin{maiorem}; it is not the noun of the
ablative standard.

\textbf{Translation.}  ``No one has love greater than this love.''
\end{workedexample}

\begin{workedexample}{ordinal time in the Creed}{W.II.7}
\Latin{Et resurrexit tertia die secundum Scripturas.}

\textbf{Analysis.}  \Latin{tertia die} is ablative of time when. Here
\Latin{dies} is feminine, as shown by feminine ablative \Latin{tertia};
\Latin{secundum} is a preposition governing the
accusative plural.

\textbf{Translation.}  ``And on the third day he rose again according to the Scriptures.''
\end{workedexample}

\begin{workedexample}{multiplicative counting (Matthew 26:34)}{W.II.8}
\Latin{Priusquam gallus cantet, ter me negabis.}

\textbf{Analysis.}  Indeclinable adverb \Latin{ter} modifies \Latin{negabis}: the denial will
occur three times.  It is not an adjective agreeing with an expressed noun.

\textbf{Translation.}  ``Before the cock crows, you will deny me three times.''
\end{workedexample}

\section{Writing laboratory}

Write the five irregular comparison rows from memory.  Under each, add one Latin phrase that
uses the comparative and one that uses the superlative.  Copy \Latin{Maiorem hac dilectionem
nemo habet} three times: first in Missal order, then with \Latin{hac dilectione} together, then
with \Latin{quam hæc dilectio}.  Explain why the case changes.

Compose: ``Christ is higher than every king'' in both comparison constructions; ``on the
second day'' and ``for three days''; ``they pray more devoutly''; ``each priest receives one
host.''  Use \Latin{altus, rex, dies, devote, singuli, sacerdos, hostia}.



\end{courseunit}

\begin{courseunit}{II.3}
\chapter{Infinitives and indirect statement}

\begin{learningobjectives}
  \item form all six ordinary infinitives and state their relative time and voice;
  \item distinguish complementary, subject, object, and indirect-statement infinitives;
  \item identify the accusative subject and predicate complement of indirect statement;
  \item preserve relative time when converting direct speech to reported speech.
\end{learningobjectives}

\section{The infinitive system}

An \Term{infinitive} is a non-finite verbal form: it expresses tense and voice
but not grammatical person. Its \Term{relative time} locates its action before,
at the same time as, or after the action of its governing verb; the label
``present infinitive'' therefore means simultaneous, not necessarily present
on the calendar. An \Term{indirect statement} reports a thought, perception,
or assertion in its ordinary form as an accusative subject plus infinitive rather than quoting it as
an independent finite clause.

\begin{grammarbox}{Forms and relative time}
\begin{center}
\begin{tabular}{@{}llll@{}}
\toprule
& Present (simultaneous) & Perfect (prior) & Future (subsequent)\\ \midrule
Active & \Latin{amare} & \Latin{amavisse} & \Latin{amaturus, -a, -um esse}\\
Passive & \Latin{amari} & \Latin{amatus, -a, -um esse} & \Latin{amatum iri}\\
\bottomrule
\end{tabular}
\end{center}

Present passive infinitives are \Latin{amari, moneri, regi, capi, audiri}: the first, second,
and fourth conjugations replace the active infinitive's final \Latin{-e} with \Latin{-i}, while
the third and third-\Latin{-io} use the present stem without its short vowel plus \Latin{-i}.
Perfect passive and future active participles agree
with their accusative subject in indirect statement: \Latin{dicit hostiam oblatam esse};
\Latin{dicit Christum venturum esse}.  The future passive \Latin{amatum iri} is invariant and
rare; ecclesiastical prose often substitutes \Latin{fore ut} + subjunctive.

The infinitive tense is relative to the governing verb, not calendar time.  Thus
\Latin{dixit eum venire} means he said the man was coming then; \Latin{dixit eum venisse}, that
he had come; \Latin{dixit eum venturum esse}, that he would come.
\end{grammarbox}

\begin{grammarbox}{Four jobs}
\begin{enumerate}
  \item A \Term{complementary infinitive} completes the meaning of verbs such
        as \Latin{possum, debeo, volo, nolo,
        incipio}: \Latin{possumus orare}.
  \item A \Term{subject infinitive} fills the subject slot and acts as a neuter
        singular noun: \Latin{orare bonum est}.
  \item Some verbs order or permit an accusative to act: \Latin{iube hæc perferri}; the
        accusative is logical subject of the passive infinitive.
  \item Verbs of saying, thinking, knowing, perceiving, and their passives introduce indirect
        statement: accusative subject + infinitive.  A reflexive \Latin{se} refers to the
        reporting subject.  Predicate nouns/adjectives are accusative: \Latin{scimus Christum
        Dominum esse}.
\end{enumerate}
\end{grammarbox}

\section{Reading the forms in context}

\begin{workedexample}{simultaneous action in reported thought (Luke 2:44)}{W.II.9}
\Latin{Existimantes autem illum esse in comitatu, venerunt iter diei.}

\textbf{Analysis.}  \Latin{Existimantes} introduces indirect statement.  Accusative
\Latin{illum} is subject of present infinitive \Latin{esse}; his supposed presence is
simultaneous with their thinking.  \Latin{Iter diei} is an accusative of extent with a
genitive measure.

\textbf{Translation.}  ``Thinking that he was in the company, they traveled a day's journey.''
\end{workedexample}

\begin{workedexample}{an ordered passive action in the Canon}{W.II.10}
\Latin{Iube hæc perferri per manus sancti Angeli tui in sublime altare tuum.}

\textbf{Analysis.}  Imperative \Latin{iube} takes accusative \Latin{hæc} as the subject of
present passive infinitive \Latin{perferri}.  The present infinitive marks action simultaneous
with the command, not present clock time.

\textbf{Translation.}  ``Command these things to be carried by the hands of your holy Angel
to your altar on high.''
\end{workedexample}

\begin{workedexample}{two dependent infinitives before Communion}{W.II.11}
\Latin{Fac me tuis semper inhærere mandatis, et a te numquam separari permittas.}

\textbf{Analysis.}  \Latin{Me} is the understood subject of active \Latin{inhærere} after
\Latin{fac}; it is also the understood subject of passive \Latin{separari} after
\Latin{permittas}.  \Latin{mandatis} is dative with \Latin{inhærere}; \Latin{a te} marks
separation.

\textbf{Translation.}  ``Make me always cling to your commandments, and never permit me to be
separated from you.''
\end{workedexample}

\begin{workedexample}{predicate accusative (Romans 6:11)}{W.II.12}
\Latin{Ita et vos existimate, vos mortuos quidem esse peccato, viventes autem Deo, in Christo
Iesu Domino nostro.}

\textbf{Analysis.}  \Latin{Existimate} introduces reported judgment.  \Latin{vos} is the
accusative subject of \Latin{esse}; predicate \Latin{mortuos} agrees with it.  Parallel
\Latin{viventes} shares the omitted \Latin{esse}.  Datives \Latin{peccato} and \Latin{Deo}
express relation.

\textbf{Translation.}  ``So you too, consider yourselves indeed dead to sin but alive to God
in Christ Jesus our Lord.''
\end{workedexample}

\section{Writing laboratory}

Copy the six-infinitive grid for \Latin{amo}, then reproduce it for \Latin{rego} and the
deponent \Latin{sequor} (present \Latin{sequi}, perfect \Latin{secutum esse}, future
\Latin{secuturum esse}).  Draw a time line under each tense.

Turn \Latin{Christus venit} (present, \English{comes}) / \Latin{venit} (perfect,
\English{came/has come}) / \Latin{veniet} (future) into statements after \Latin{Sacerdos dicit} and
after \Latin{Sacerdos dixit}.  Compose ``We believe Christ to be Lord,'' ``He said that the
offering had been received,'' and ``Command the gifts to be blessed.''  Explicitly mark each
accusative subject.



\end{courseunit}

\begin{courseunit}{II.4}
\chapter{Participles and the ablative absolute}

\begin{learningobjectives}
  \item form and translate present active, perfect passive, and future active participles;
  \item handle active-perfect participles of deponents and the irregular \Latin{mortuus};
  \item distinguish attributive, circumstantial, supplementary, and periphrastic uses;
  \item recognize and unpack an ablative absolute without inventing a false subject.
\end{learningobjectives}

\section{Forms and logic}

A \Term{participle} is a verbal adjective: it agrees with a noun in gender,
number, and case while retaining verbal features such as voice, relative time,
objects, and adverbial modifiers. An \Term{attributive participle} directly
modifies or identifies a noun. A \Term{circumstantial participle} supplies the
time, cause, concession, means, or attendant circumstance of an action. A
\Term{supplementary participle} completes the sense of certain verbs of
perceiving or depicting. A \Term{periphrastic} combines a participle or
gerundive with a form of \Latin{sum} to express a verbal idea as a unit.

\begin{grammarbox}{The participial system}
\begin{tabularx}{\linewidth}{@{}P{1.25in}P{1.55in}P{1.15in}Y@{}}
\toprule
\textbf{Participle} & \textbf{Model} & \textbf{Relative time} & \textbf{Voice and formation}\\ \midrule
Present active & \Latin{amans, amantis} & simultaneous & active; present stem +
\Latin{-ns/-ntis}, third-declension one-termination adjective\\
Perfect passive & \Latin{amatus, -a, -um} & prior & passive; fourth principal part,
first/second-declension adjective\\
Future active & \Latin{amaturus, -a, -um} & subsequent or intended & active; supine stem +
\Latin{-urus}\\
Gerundive & \Latin{amandus, -a, -um} & no tense & passive necessity or purpose; treated fully
in Unit 12\\
\bottomrule
\end{tabularx}

Latin lacks an ordinary perfect active participle, but deponents supply one with active
meaning: \Latin{secutus}, ``having followed''; \Latin{locutus}, ``having spoken.''  Likewise
\Latin{mortuus} means ``having died/dead.''  Participles agree with a noun in gender, number,
and case while retaining verbal powers such as objects and adverbs.
\end{grammarbox}

\begin{grammarbox}{Ablative absolute}
An ablative absolute normally contains a noun/pronoun plus a participle, both ablative, and is
grammatically detached (\Latin{absolutus}) from the main clause.  Translate the logical
relation demanded by context: ``when,'' ``after,'' ``because,'' ``although,'' or
``with \ldots{}.''  Present participle = action simultaneous with the main verb; perfect
participle = prior action.  A noun or adjective may stand without a participle because Latin
has no present participle of \Latin{sum}: \Latin{Christo duce}, ``with Christ as leader.''

Do not call a phrase absolute if its noun is also a required argument of the main verb.  Do not
automatically translate a perfect passive participle actively: \Latin{pane accepto} is ``when
the bread had been received/taken,'' with context supplying the agent.
\end{grammarbox}

\section{Reading the forms in context}

\begin{workedexample}{present participle in the institution narrative}{W.II.13}
\Latin{Accipiens et hunc præclarum calicem in sanctas ac venerabiles manus suas, item tibi
gratias agens, benedixit.}

\textbf{Analysis.}  Nominative \Latin{accipiens} and \Latin{agens} agree with the implicit
subject Christ and denote actions simultaneous with or immediately attendant on
\Latin{benedixit}.  \Latin{calicem} is object of \Latin{accipiens}; \Latin{tibi} is dative
with \Latin{gratias agens}.

\textbf{Translation.}  ``Taking also this noble chalice into his holy and venerable hands,
and likewise giving thanks to you, he blessed it.''
\end{workedexample}

\begin{workedexample}{perfect participles in the Creed}{W.II.14}
\Latin{Genitum, non factum, consubstantialem Patri.}

\textbf{Analysis.}  Accusative masculine singular \Latin{genitum} and \Latin{factum} agree
with \Latin{Filium} from the preceding line; both are perfect passive participles.  The
contrast is theological and grammatical: ``begotten,'' not ``made.''  Predicate adjective
\Latin{consubstantialem} also agrees with \Latin{Filium}.

\textbf{Translation.}  ``Begotten, not made, of one substance with the Father.''
\end{workedexample}

\begin{workedexample}{ablative absolute in the Canon}{W.II.15}
\Latin{Qui, pridie quam pateretur, accepit panem in sanctas ac venerabiles manus suas, et
elevatis oculis in cælum \ldots{} benedixit.}

\textbf{Analysis.}  In \Latin{elevatis oculis}, both words are ablative masculine plural.
The perfect passive participle marks an action prior to \Latin{benedixit}; \Latin{oculis} has
no syntactic role in the main clause.

\textbf{Translation.}  ``He, on the day before he suffered, took bread into his holy and
venerable hands and, when his eyes had been raised to heaven, blessed it.''
\end{workedexample}

\begin{workedexample}{participles sharing the main subject (Matthew 2:11)}{W.II.16}
\Latin{Et intrantes domum invenerunt puerum cum Maria matre eius, et procidentes adoraverunt
eum.}

\textbf{Analysis.}  \Latin{Intrantes} and \Latin{procidentes} are present active nominative
plural participles, from active \Latin{intro} and \Latin{procido}; they introduce attendant
actions of the same people who \Latin{invenerunt} and \Latin{adoraverunt}.  Neither is an
ablative absolute because both share the main subject.  Compare Matthew 8:23,
\Latin{secuti sunt eum discipuli eius}: deponent \Latin{secuti sunt} has perfect morphology
but active meaning, ``his disciples followed him.''

\textbf{Translation.}  ``Entering the house, they found the child with Mary his mother, and
falling down they adored him.''
\end{workedexample}

\section{Writing laboratory}

Write all four participial stems for \Latin{amo, moneo, rego, capio, audio}; decline
\Latin{regens} in singular and plural and \Latin{benedictus} in all genders.  Copy
\Latin{elevatis oculis in cælum} and draw a box around the absolute phrase.

Compress these pairs: \Latin{Sacerdos orat et hostiam offert} $\rightarrow$
\Latin{Sacerdos orans hostiam offert}; \Latin{Oculi elevati sunt; Christus benedixit}
$\rightarrow$ \Latin{Oculis elevatis, Christus benedixit}.  Compose one ablative absolute
with a present participle and one with a perfect participle, then expand each into a finite
\Latin{cum}-clause.



\end{courseunit}

\begin{courseunit}{II.5}
\chapter{Subjunctive morphology and sequence of tenses}

\begin{learningobjectives}
  \item form every active and passive subjunctive of a regular verb and \Latin{sum};
  \item explain the time relation expressed by each subjunctive tense;
  \item apply primary and historic sequence in dependent clauses;
  \item distinguish a morphological identification from a clause's syntactic use.
\end{learningobjectives}

\section{The four tenses}

The \Term{subjunctive mood} presents an action through the speaker's or
sentence's framing---for example as wished, commanded, purposed, resulting,
possible, or dependent in reported thought. The \Term{indicative} ordinarily
asserts or asks about an action as a fact. Mood is a property of the verb form;
the clause type explains why that mood is used in context.

When a subjunctive stands independently, a \Term{jussive} expresses a
third-person command, a \Term{hortatory} urges first-person plural action, and
an \Term{optative} expresses a wish. In dependent clauses, the conjunction and
the clause's function determine its use. These independent uses return in the
unit on commands and wishes.

\begin{grammarbox}{Present and imperfect subjunctive}
The present subjunctive uses the tense signs \Latin{-e-} in the first conjugation and
\Latin{-a-} in the others, followed by personal endings:

\begin{center}
\begin{tabular}{@{}llllll@{}}
\toprule
\Latin{amare} & \Latin{monere} & \Latin{regere} & \Latin{capere} & \Latin{audire} &
\Latin{esse}\\ \midrule
amem & moneam & regam & capiam & audiam & sim\\
ames & moneas & regas & capias & audias & sis\\
amet & moneat & regat & capiat & audiat & sit\\
amemus & moneamus & regamus & capiamus & audiamus & simus\\
ametis & moneatis & regatis & capiatis & audiatis & sitis\\
ament & moneant & regant & capiant & audiant & sint\\
\bottomrule
\end{tabular}
\end{center}

The imperfect adds personal endings directly to the present active infinitive:
\Latin{amarem, amares, amaret, amaremus, amaretis, amarent}; \Latin{essem, esses, esset,
essemus, essetis, essent}.  Passive forms replace active endings with
\Latin{-r, -ris/-re, -tur, -mur, -mini, -ntur}: \Latin{amer, ameris, ametur};
\Latin{amarer, amareris, amaretur}.  Deponents use passive morphology with active meaning.
\end{grammarbox}

\begin{grammarbox}{Perfect and pluperfect subjunctive}
The perfect active is perfect stem + \Latin{-eri-} + personal endings:
\Latin{amaverim, amaveris, amaverit, amaverimus, amaveritis, amaverint}.  The pluperfect
active is the perfect active infinitive + personal endings:
\Latin{amavissem, amavisses, amavisset, amavissemus, amavissetis, amavissent}.

Perfect passive is perfect passive participle + present subjunctive of \Latin{sum}:
\Latin{amatus sim}; pluperfect passive uses imperfect subjunctive:
\Latin{amatus essem}.  The participle always agrees with the subject:
\Latin{hostiæ oblatæ sint}; \Latin{dona benedicta essent}.

Beware identical-looking forms.  \Latin{amaverit} can be perfect subjunctive or future perfect
indicative; \Latin{amaverimus} can be perfect subjunctive or future perfect indicative.
Clause environment, not appearance alone, decides.
\end{grammarbox}

\section{Sequence of tenses}

\Term{Sequence of tenses} is the conventional matching of a dependent
subjunctive's tense to the time of its governing verb and to whether the
dependent action is simultaneous/subsequent or prior. \Term{Primary sequence}
uses a present or future viewpoint; \Term{historic sequence} uses a past
viewpoint. Sequence locates the dependent action; it does not name the clause's
purpose or meaning.

\begin{grammarbox}{The sequence table}
\begin{center}
\begin{tabularx}{0.92\linewidth}{@{}P{1.4in}Y Y@{}}
\toprule
\textbf{Governing verb} & \textbf{Same or subsequent dependent action} &
\textbf{Prior dependent action}\\ \midrule
Primary: present, future, future perfect, or perfect with present force & present subjunctive &
perfect subjunctive\\
Historic: imperfect, historical perfect, or pluperfect & imperfect subjunctive & pluperfect
subjunctive\\
\bottomrule
\end{tabularx}
\end{center}

Sequence locates dependent action relative to the governing action.  It does not itself tell
why the subjunctive is used.  First label the clause---purpose, result, indirect command,
question, \Latin{cum}, and so forth---then state sequence.  Writers may use primary sequence
after a historical present, preserve a perfect for a completed result, or use a tense required
by an independent condition.  Those are meaningful choices, not permission to ignore the
table.
\end{grammarbox}

\section{Reading the forms in context}

\begin{workedexample}{present sequence at the Orate fratres}{W.II.17}
\Latin{Orate, fratres, ut meum ac vestrum sacrificium acceptabile fiat apud Deum Patrem
omnipotentem.}

\textbf{Analysis.}  \Latin{Fiat} is third-person singular present subjunctive of \Latin{fio};
it follows the present-time imperative \Latin{orate} and expresses action subsequent to the
prayer.  Morphology is present; the English sense is ``may become/be made.''

\textbf{Translation.}  ``Pray, brethren, that my sacrifice and yours may be acceptable before
God, the almighty Father.''
\end{workedexample}

\begin{workedexample}{anterior action within the Canon}{W.II.18}
\Latin{Ut quotquot ex hac altaris participatione sacrosanctum Filii tui Corpus et Sanguinem
sumpserimus, omni benedictione cælesti et gratia repleamur.}

\textbf{Analysis.}  Perfect subjunctive \Latin{sumpserimus} is prior to present subjunctive
\Latin{repleamur}.  The latter is first-person plural present passive.  The sequence asks that
all of us who have received may then be filled.

\textbf{Translation.}  ``That all of us who have received the most holy Body and Blood of your
Son from this sharing in the altar may be filled with every heavenly blessing and grace.''
\end{workedexample}

\begin{workedexample}{historic sequence before the Passion}{W.II.19}
\Latin{Qui pridie quam pateretur, accepit panem.}

\textbf{Analysis.}  \Latin{Pateretur} is third-person singular imperfect subjunctive of the
deponent \Latin{patior}.  With historic \Latin{accepit}, \Latin{pridie quam} places the
anticipated suffering after the taking of bread.

\textbf{Translation.}  ``Who, on the day before he was to suffer, took bread.''
\end{workedexample}

\begin{workedexample}{historic indirect command (Mark 5:10)}{W.II.20}
\Latin{Et deprecabatur eum multum, ne se expelleret extra regionem.}

\textbf{Analysis.}  Imperfect \Latin{deprecabatur} governs an indirect command.  Imperfect
subjunctive \Latin{expelleret} is same/subsequent in historic sequence; \Latin{se} is its
accusative object and refers to the speaker.

\textbf{Translation.}  ``And he earnestly begged him not to drive him out of the region.''
\end{workedexample}

\section{Writing laboratory}

Write from memory the present active table for all five conjugations and \Latin{sum}.  Beneath
each form write its passive counterpart.  Then produce all four subjunctive tenses of
\Latin{benedico} in third-person singular and plural, active and passive.

Build two time diagrams.  Under \Latin{rogat}, place \Latin{veniat} and \Latin{venerit}; under
\Latin{rogavit}, place \Latin{veniret} and \Latin{venisset}.  Compose one Latin sentence for
each diagram and explain its relative time in English.



\end{courseunit}

\begin{courseunit}{II.6}
\chapter{Purpose, result, and indirect command}

\begin{learningobjectives}
  \item distinguish three outwardly similar \Latin{ut}-clause families by meaning and signal;
  \item use \Latin{ne} for negative purpose or command but \Latin{ut non} for negative result;
  \item recognize relative purpose;
  \item preserve sequence when transforming phrases into clauses.
\end{learningobjectives}

\section{Three clause patterns}

A \Term{purpose clause} states the intended end of an action. A \Term{result
clause} states the consequence that follows from a degree or circumstance. An
\Term{indirect command} reports what someone asks, orders, urges, permits, or
prevents another person from doing. Because all three may contain \Latin{ut}
and a subjunctive, identify the clause's job before translating the marker.

\begin{grammarbox}{Purpose}
Purpose answers ``for what end?'' and uses \Latin{ut} + subjunctive, negative \Latin{ne}:
\Latin{venit ut oret}; \Latin{tacet ne turbet}.  A relative pronoun may replace
\Latin{ut is}: \Latin{misit sacerdotem qui oraret}, ``he sent a priest to pray.''  The relative
then has a case-role within its own clause.  Purpose is necessarily simultaneous with or
subsequent to the governing action, so primary sequence uses present and historic sequence
uses imperfect.
\end{grammarbox}

\begin{grammarbox}{Result}
Result states what actually or characteristically follows.  It uses \Latin{ut}; its negative
is \Latin{ut non}, \Latin{ut nemo}, or \Latin{ut nihil}---not purpose \Latin{ne}.  Typical
signals are \Latin{tam, ita, sic, tantus, talis, tot, adeo} or a comparative.  Perfect
subjunctive may emphasize a completed result even after a historic verb.
\end{grammarbox}

\begin{grammarbox}{Indirect command}
Verbs of asking, commanding, urging, permitting, persuading, and taking care introduce the
content of a command with \Latin{ut/ne} + subjunctive: \Latin{rogat ut veniat};
\Latin{præcipit ne discedant}.  Ask whether the clause tells what someone is requested to do
(indirect command) or why an action is done (purpose).  Some command verbs instead take an
accusative plus infinitive: \Latin{iube eum venire}; do not force both constructions together.
\end{grammarbox}

\section{Reading the forms in context}

\begin{workedexample}{purpose (John 1:7)}{W.II.21}
\Latin{Hic venit in testimonium, ut testimonium perhiberet de lumine, ut omnes crederent per
illum.}

\textbf{Analysis.}  Both \Latin{ut}-clauses express purpose after historic \Latin{venit}.
Imperfect subjunctives \Latin{perhiberet} and \Latin{crederent} follow historic sequence;
\Latin{per illum} identifies the human instrument.

\textbf{Translation.}  ``He came as a witness, to bear witness concerning the light, so that
all might believe through him.''
\end{workedexample}

\begin{workedexample}{negative purpose (Matthew 26:41)}{W.II.22}
\Latin{Vigilate et orate, ut non intretis in tentationem.}

\textbf{Analysis.}  The Vulgate here has \Latin{ut non}; context makes the clause the intended
end of watching and praying.  It therefore shows that the semantic test outranks a mechanical
claim that every \Latin{ut non} must be result.  Present \Latin{intretis} follows imperatives.

\textbf{Translation.}  ``Keep watch and pray, so that you may not enter into temptation.''
\end{workedexample}

\begin{workedexample}{result (John 3:16)}{W.II.23}
\Latin{Sic enim Deus dilexit mundum, ut Filium suum unigenitum daret.}

\textbf{Analysis.}  Correlative \Latin{sic \ldots{} ut} measures the extent and consequence of
God's love.  Imperfect subjunctive \Latin{daret} follows historic \Latin{dilexit}.

\textbf{Translation.}  ``For God so loved the world that he gave his only-begotten Son.''
\end{workedexample}

\begin{workedexample}{indirect command (Matthew 14:28)}{W.II.24}
\Latin{Domine, si tu es, iube me ad te venire super aquas.}

\textbf{Analysis.}  \Latin{Iube} selects the alternative accusative-plus-infinitive command:
\Latin{me} is subject of \Latin{venire}.  Recasting with a verb such as \Latin{præcipe}
would give \Latin{præcipe ut ad te veniam}.  The purpose is to recognize, not blend, the two
patterns.

\textbf{Translation.}  ``Lord, if it is you, command me to come to you over the waters.''
\end{workedexample}

\begin{workedexample}{relative purpose (Numbers 31:4)}{W.II.25}
\Latin{Mille viri de singulis tribubus eligantur ex Israël qui mittantur ad bellum.}

\textbf{Analysis.}  Jussive \Latin{eligantur} orders the selection.  Relative \Latin{qui} is
nominative masculine plural, referring to \Latin{viri}, and present passive subjunctive
\Latin{mittantur} states the purpose for which they are chosen.

\textbf{Translation.}  ``Let a thousand men from each tribe be chosen out of Israel to be sent
to war.''
\end{workedexample}

\section{Writing laboratory}

Copy the headings \textsc{purpose / result / indirect command}.  Under each, write its
question, positive marker, negative marker, and two typical signals.  Then bracket and label
every clause in the five worked examples.

Transform \Latin{sacerdos venit; sacrificium offert} into a purpose clause; expand
\Latin{tam sanctus est} with a result; report \Latin{Venite ad altare} after \Latin{Sacerdos
monuit fratres}.  Make each negative, paying special attention to \Latin{ne} versus
\Latin{ut non}.



\end{courseunit}

\begin{courseunit}{II.7}
\chapter{Substantive clauses: a clause acting as a noun}

\begin{learningobjectives}
  \item identify a clause serving as subject, object, apposition, or complement;
  \item distinguish volitive \Latin{ut/ne}, fact/result \Latin{ut/ut non}, and factual
        \Latin{quod/quia/quoniam};
  \item recognize \Latin{quin} and \Latin{quominus} after negative hindrance or doubt;
  \item form indirect questions with an interrogative word, subjunctive, and proper sequence;
  \item analyze the conventional request structure of collects and the Canon.
\end{learningobjectives}

\section{Clause families}

\begin{grammarbox}{Volitive and result/content clauses}
A \Term{substantive clause} fills a noun slot as subject, object, apposition, or
predicate complement; it does not merely modify a verb adverbially. A
\Term{volitive clause} communicates an act of will. Volitive clauses follow
verbs and expressions of wishing, asking, causing, permitting, or
deciding and use \Latin{ut/ne} + subjunctive.  Indirect command is the most concrete subtype.
A \Term{result/content clause} states the fact or outcome that completes an
impersonal expression. Such clauses use \Latin{ut/ut non} + subjunctive and
occur after \Latin{fit, accidit, efficitur, restat} and expressions such as
\Latin{dignum est}. No degree word is required
when the clause itself completes the idea.

The question test is useful: in \Latin{rogamus ut venias}, ``what do we ask?''---the clause is
object/content.  In \Latin{venimus ut oremus}, ``why do we come?''---the clause is purpose.
\end{grammarbox}

\begin{grammarbox}{Facts, doubt, and prevention}
\Latin{Quod, quia}, and \Latin{quoniam} + indicative may introduce a fact (``the fact that'')
or a stated reason.  A subjunctive can distance the author from a reason attributed to someone
else.  After a negative expression of doubt, \Latin{quin} + subjunctive means ``that'':
\Latin{non dubito quin audiat}, ``I do not doubt that he hears.''  After prevention or refusal,
\Latin{quominus} or \Latin{ne} + subjunctive means ``from doing'':
\Latin{nihil impedit quominus oremus}.  These last patterns are less frequent in the Ordinary
but prevent misreading in lessons and rubrics.
\end{grammarbox}

\begin{grammarbox}{Indirect questions}
A direct question becomes a substantive clause after asking, knowing, seeing, considering, or
doubting: \Latin{Quid facit?} becomes \Latin{scio quid faciat}; after a historic verb,
\Latin{sciebam quid faceret}.  The interrogative keeps the case demanded inside its clause,
but the finite verb is subjunctive and follows sequence.  Yes/no questions use \Latin{num};
alternatives use \Latin{utrum \ldots{} an} or simply \Latin{-ne \ldots{} an}.  Do not confuse
interrogative \Latin{quid faciat}, ``what he does,'' with headless relative
\Latin{quod facit}, ``what/the thing that he does.''
\end{grammarbox}

\begin{grammarbox}{The liturgical request frame}
A collect often has this architecture:
\begin{quote}
\Latin{Deus} [address], \Latin{qui \ldots{}} [relative ground],
\Latin{concede/quæsumus} [request], \Latin{ut \ldots{}} [content or desired result].
\end{quote}
The colon or semicolon does not create the syntax.  Find the finite request, then attach the
\Latin{ut}-clause as what is requested.  Vocatives, appositions, and a long relative clause
may stand between them.
\end{grammarbox}

\section{Reading the forms in context}

\begin{workedexample}{the Canon's substantive request}{W.II.26}
\Latin{Te igitur, clementissime Pater, supplices rogamus ac petimus, uti accepta habeas et
benedicas hæc dona, hæc munera, hæc sancta sacrificia illibata.}

\textbf{Analysis.}  Archaic \Latin{uti} equals \Latin{ut}.  Its two present subjunctives
\Latin{habeas} and \Latin{benedicas} give the content of \Latin{rogamus ac petimus}.  Fronted
\Latin{te} is object of the asking verbs; \Latin{supplices} agrees with the implicit
\Latin{nos}.

\textbf{Translation.}  ``Therefore, most merciful Father, we humbly ask and petition you to
hold acceptable and bless these gifts, these offerings, these holy and unblemished
sacrifices.''
\end{workedexample}

\begin{workedexample}{content after worthiness before Communion}{W.II.27}
\Latin{Domine, non sum dignus ut intres sub tectum meum.}

\textbf{Analysis.}  \Latin{Ut intres} completes the respect in which the speaker is not
worthy; it is a substantive clause of content, not the purpose for being unworthy.  Present
subjunctive \Latin{intres} follows present \Latin{sum}.

\textbf{Translation.}  ``Lord, I am not worthy that you should enter beneath my roof.''
\end{workedexample}

\begin{workedexample}{subject/content after happening (Luke 10:31)}{W.II.28}
\Latin{Accidit autem ut sacerdos quidam descenderet eadem via.}

\textbf{Analysis.}  The \Latin{ut}-clause supplies what happened and is the logical subject of
impersonal \Latin{accidit}.  Imperfect \Latin{descenderet} follows the historical perfect.

\textbf{Translation.}  ``It happened that a certain priest was going down by the same road.''
\end{workedexample}

\begin{workedexample}{a stated fact (Galatians 4:6)}{W.II.29}
\Latin{Quoniam autem estis filii, misit Deus Spiritum Filii sui in corda vestra.}

\textbf{Analysis.}  \Latin{Quoniam} introduces as fact the ground for \Latin{misit};
\Latin{estis} is indicative, not subjunctive.  The clause is causal in relation to the main
sentence, while its proposition is asserted as true.

\textbf{Translation.}  ``And because you are sons, God sent the Spirit of his Son into your
hearts.''
\end{workedexample}

\begin{workedexample}{an indirect question (John 6:6)}{W.II.30}
\Latin{Hoc autem dicebat tentans eum: ipse enim sciebat quid esset facturus.}

\textbf{Analysis.}  Interrogative \Latin{quid} is the neuter accusative object of the action
in \Latin{facturus}.  Imperfect subjunctive \Latin{esset}, governed by historic
\Latin{sciebat}, combines with the future active participle to make that action subsequent to
the knowing.

\textbf{Translation.}  ``He said this to test him, for he himself knew what he was going to
do.''
\end{workedexample}

\section{Writing laboratory}

Copy the following Collect printed for this exercise: \Latin{Excita, quæsumus,
Domine, potentiam tuam, et veni: ut ab imminentibus peccatorum nostrorum
periculis te mereamur protegente eripi, te liberante salvari.} Box the address,
underline the two commands once, and bracket the substantive content governed
by \Latin{mereamur}. Rewrite the sentence in diagnostic prose order without
deleting a word. (\emph{Source:} \Missal, Dominica I Adventus, Collecta.)

Compose: ``We ask that you bless these gifts''; ``We ask that you not abandon us''; ``It
happened that the priest came''; ``I do not doubt that God hears''; ``Nothing prevents us from
praying.''  Use \Latin{rogare, benedicere, derelinquere, accidere, dubitare, impedire} and
state why each marker is \Latin{ut, ne, quin}, or \Latin{quominus}.



\end{courseunit}

\begin{courseunit}{II.8}
\chapter{Temporal, causal, and concessive clauses}

\begin{learningobjectives}
  \item choose indicative or subjunctive with \Latin{cum} according to meaning;
  \item read the principal temporal conjunctions and their time relations;
  \item distinguish asserted cause from circumstantial or reported cause;
  \item recognize concession with \Latin{quamquam, etsi, licet}, and \Latin{cum}.
\end{learningobjectives}

\section{Mood follows viewpoint}

A \Term{temporal clause} locates one action in relation to another. A
\Term{causal clause} supplies a reason. A \Term{concessive clause} admits a
circumstance despite which the main statement remains true (``although''). The
same conjunction can introduce more than one relation, so meaning and mood
must be considered together.

\begin{grammarbox}{\Latin{Cum}}
\begin{tabularx}{\linewidth}{@{}P{1.7in}P{1.0in}Y@{}}
\toprule
\textbf{Use} & \textbf{Mood} & \textbf{Test}\\ \midrule
temporal, definite & indicative & ``when'' identifies a time as a fact, especially present or
repeated action\\
circumstantial/historical & subjunctive & supplies the background in which the main event
unfolds; usually imperfect or pluperfect in past narration\\
causal & subjunctive & ``since/because,'' presenting a reason as the writer's framing\\
concessive & subjunctive & ``although''; main clause often has \Latin{tamen}\\
\bottomrule
\end{tabularx}

With indicative, \Latin{cum} may be temporal or causal when the cause is asserted as an
objective fact.  Translate only after testing the relation.
\end{grammarbox}

\begin{grammarbox}{Other temporal and causal markers}
\Latin{Ubi, ut, postquam} normally take indicative for a completed factual event.
\Latin{Dum} + present indicative commonly means ``while'' even beside a past main verb;
\Latin{dum} may also take subjunctive for expectation or purpose.  \Latin{donec, quoad}, and
\Latin{antequam/priusquam} take indicative for an actual fact and subjunctive when the event
is anticipated, intended, or merely possible.  \Latin{Simul ac} means ``as soon as.''

\Latin{Quia, quoniam, quod} + indicative assert a cause; subjunctive can mark an alleged or
reported cause.  \Latin{Quamquam/etsi} generally take indicative for admitted fact;
\Latin{licet} and concessive \Latin{cum} take subjunctive.  \Latin{Etiamsi} follows the mood of
the condition it introduces.
\end{grammarbox}

\section{Reading the forms in context}

\begin{workedexample}{historical \Latin{cum} (Matthew 2:1)}{W.II.31}
\Latin{Cum ergo natus esset Iesus in Bethlehem Iuda in diebus Herodis regis, ecce Magi ab
oriente venerunt Ierosolymam.}

\textbf{Analysis.}  Pluperfect subjunctive \Latin{natus esset} gives prior historical
background to \Latin{venerunt}; it is not a contrary-to-fact condition.  The clause means
``after/when Jesus had been born.''

\textbf{Translation.}  ``When Jesus had therefore been born in Bethlehem of Judah in the days
of King Herod, behold, Magi came from the east to Jerusalem.''
\end{workedexample}

\begin{workedexample}{two \Latin{dum}-clauses in a Christmas Introit}{W.II.32}
\Latin{Dum medium silentium tenerent omnia, et nox in suo cursu medium iter haberet,
omnipotens sermo tuus, Domine, de cælis a regalibus sedibus venit.}

\textbf{Analysis.}  Imperfect subjunctives \Latin{tenerent} and \Latin{haberet} frame the
circumstance in which \Latin{venit} occurs.  \Latin{Omnia} is the plural subject of
\Latin{tenerent}; \Latin{sermo} is subject of \Latin{venit}.

\textbf{Translation.} ``While all things held the middle silence and night
held the middle of its course, your almighty Word, O Lord, came from heaven, from royal
thrones.''
\end{workedexample}

\begin{workedexample}{anticipated time (Matthew 26:34)}{W.II.33}
\Latin{Priusquam gallus cantet, ter me negabis.}

\textbf{Analysis.}  The cock's crow is anticipated relative to future \Latin{negabis}; hence
present subjunctive \Latin{cantet}.  \Latin{Priusquam} is temporal, not a purpose marker.

\textbf{Translation.}  ``Before the cock crows, you will deny me three times.''
\end{workedexample}

\begin{workedexample}{concessive circumstance (2 Corinthians 8:9)}{W.II.34}
\Latin{Propter vos egenus factus est, cum esset dives.}

\textbf{Analysis.}  \Latin{Cum esset} cannot mean simple chronological sequence here; the
contrast with \Latin{egenus factus est} makes it concessive: ``although he was rich.''

\textbf{Translation.}  ``For your sake he became poor, although he was rich.''
\end{workedexample}

\Needspace{20\baselineskip}
\section{Writing laboratory}

Make a four-column page headed \textsc{definite time / background / cause / concession}.
Copy each worked clause into its column and write mood, tense, and viewpoint beneath it.

Join \Latin{Christus natus est; Magi venerunt} with historical \Latin{cum}; join
\Latin{gallus cantabit; Petrus negabit} with anticipated \Latin{priusquam}; concede
\Latin{dives erat} before \Latin{egenus factus est}; and assert a cause with \Latin{quia}.
Then change the asserted cause into a reported allegation by changing mood and explaining the
new stance.



\end{courseunit}

\begin{courseunit}{II.9}
\chapter{Conditions}

\begin{learningobjectives}
  \item classify open, future, and contrary-to-fact conditions;
  \item choose tense and mood independently in protasis and apodosis;
  \item recognize \Latin{nisi}, \Latin{sin}, and conditional relative words;
  \item translate conditions without turning possibility into fact or unreality into doubt.
\end{learningobjectives}

\section{The conditional system}

\begin{grammarbox}{Protasis and apodosis}
The \Term{protasis} is the condition, normally introduced by \Latin{si} or \Latin{nisi}; the
\Term{apodosis} is its consequence.

\begin{tabularx}{\linewidth}{@{}P{1.45in}P{1.55in}P{1.55in}Y@{}}
\toprule
\textbf{Type} & \textbf{Protasis} & \textbf{Apodosis} & \textbf{Normal sense}\\ \midrule
open, present/past & indicative & any appropriate indicative/imperative & leaves truth open:
``if it is/was''\\
future more vivid & future or future perfect indicative & future indicative or imperative &
expected possibility: ``if it happens''\\
future less vivid & present subjunctive & present subjunctive & tentative possibility:
``if it should happen, it would''\\
present contrary to fact & imperfect subjunctive & imperfect subjunctive & ``if it were
(but is not), it would''\\
past contrary to fact & pluperfect subjunctive & pluperfect subjunctive & ``if it had
(but did not), it would have''\\
\bottomrule
\end{tabularx}

Mixed conditions use the tense appropriate to each time sphere.  An imperative, jussive, or
question may stand in the apodosis.  \Latin{Nisi} is ``unless/if not''; \Latin{si non} usually
negates a particular word; \Latin{sin} means ``but if.''  \Latin{Quicumque} and
\Latin{si quis} can generalize a condition.
\end{grammarbox}

\begin{grammarbox}{Do not diagnose by one form}
The ending \Latin{-erit} may be future perfect indicative or perfect subjunctive.  In a vivid
future protasis paired with a future apodosis it is normally future perfect:
\Latin{si venerit, gaudebimus}.  Imperfect and pluperfect subjunctives do not always signal a
condition; conjunction and clause relation must also be present.
\end{grammarbox}

\section{Reading the forms in context}

\begin{workedexample}{vivid future (Psalm 129:3)}{W.II.35}
\Latin{Si iniquitates observaveris, Domine, Domine, quis sustinebit?}

\textbf{Analysis.}  In this future frame, \Latin{observaveris} is second-person singular
future perfect indicative, prior to future \Latin{sustinebit}.  The rhetorical question forms
the apodosis.

\textbf{Translation.}  ``If you keep watch over iniquities, O Lord, O Lord, who will endure?''
\end{workedexample}

\begin{workedexample}{open condition with jussive consequence (Matthew 16:24)}{W.II.36}
\Latin{Si quis vult post me venire, abneget semetipsum, et tollat crucem suam, et sequatur me.}

\textbf{Analysis.}  Present indicative \Latin{vult} leaves the condition open.  Three present
subjunctives---\Latin{abneget, tollat, sequatur}---are jussive apodoses, not a future-less-vivid
pair.

\textbf{Translation.}  ``If anyone wishes to come after me, let him deny himself, take up his
cross, and follow me.''
\end{workedexample}

\begin{workedexample}{future condition (John 6:54)}{W.II.37}
\Latin{Nisi manducaveritis carnem Filii hominis et biberitis eius sanguinem, non habebitis
vitam in vobis.}

\textbf{Analysis.}  \Latin{Manducaveritis} and \Latin{biberitis} are future perfect
indicatives in the protasis; \Latin{habebitis} is future indicative.  \Latin{Nisi} negates the
condition as a whole.

\textbf{Translation.}  ``Unless you eat the flesh of the Son of Man and drink his blood, you
will not have life in you.''
\end{workedexample}

\begin{workedexample}{past contrary to fact (1 Corinthians 2:8)}{W.II.38}
\Latin{Si enim cognovissent, numquam Dominum gloriæ crucifixissent.}

\textbf{Analysis.}  Both \Latin{cognovissent} and \Latin{crucifixissent} are pluperfect
subjunctive.  The construction denies the past premise and consequence.

\textbf{Translation.}  ``For if they had known, they would never have crucified the Lord of
glory.''
\end{workedexample}

\begin{workedexample}{a personal past contrary condition (John 11:21)}{W.II.39}
\Latin{Domine, si fuisses hic, frater meus non fuisset mortuus.}

\textbf{Analysis.}  Pluperfect subjunctives \Latin{fuisses} and \Latin{fuisset} place both
unfulfilled facts in the past.  \Latin{Mortuus} is a perfect participle agreeing with
\Latin{frater}.

\textbf{Translation.}  ``Lord, if you had been here, my brother would not have died.''
\end{workedexample}

\section{Writing laboratory}

Copy the five-row conditional table.  For each row write one paired verb form of \Latin{venio}
and \Latin{gaudeo}.  Mark indicative forms with a straight underline and subjunctives with a
wavy underline.

Transform \Latin{si venit, gaudemus} successively into future more vivid, future less vivid,
present contrary, and past contrary.  Then compose a condition whose apodosis is an imperative
and one introduced by \Latin{nisi}.  Translate each so its degree of reality is audible.



\end{courseunit}

\begin{courseunit}{II.10}
\chapter{Relative clauses of characteristic and related uses}

\begin{learningobjectives}
  \item distinguish a factual relative clause from one characterizing a kind of person or
        thing;
  \item recognize characteristic after negative, indefinite, or class-defining antecedents;
  \item separate characteristic from relative purpose, cause, and concession;
  \item account for mood rather than translating every relative identically.
\end{learningobjectives}

\section{Fact or characteristic?}

A \Term{factual relative clause} asserts information about an antecedent
treated as definite. A \Term{relative clause of characteristic} describes the kind or
class of person or thing represented by its antecedent rather than asserting a
single identifying fact. The difference is one of viewpoint, not of the
relative pronoun's form.

\begin{grammarbox}{Indicative and subjunctive relatives}
An indicative relative normally adds an asserted fact about a definite or treated-as-definite
antecedent: \Latin{agnus qui tollit}, ``the Lamb who takes away.''  A subjunctive relative
characterizes the sort of antecedent meant: \Latin{nemo est qui faciat}, ``there is no one of
the sort who does.''  Common frames are \Latin{sunt qui} (there are some who),
\Latin{nemo/nullus est qui}, \Latin{quis est qui?}, and a general class after
\Latin{dignus, indignus, idoneus} or a superlative/restriction.

The rule is semantic, not automatic.  \Latin{Quis est qui vincit?} may take indicative when
the question asks the identity of an actual victor.  Conversely, a definite-looking
antecedent can take subjunctive if the writer presents it as a type.
\end{grammarbox}

\begin{grammarbox}{Other subjunctive relatives}
A relative clause can express purpose (``to do''), cause (``since he''), or concession
(``although he'').  Negative relative purpose uses \Latin{qui non}, not \Latin{ne} alone.
Characteristic follows sequence of tenses when dependent on a governing past verb, but a
timeless generalization may retain present subjunctive.  Always give two labels: the relative
word's case-role and the clause's semantic use.
\end{grammarbox}

\section{Reading the forms in context}

\begin{workedexample}{negative antecedent (Psalm 13:3)}{W.II.40}
\Latin{Non est qui faciat bonum, non est usque ad unum.}

\textbf{Analysis.}  \Latin{Qui} is nominative singular and subject of present subjunctive
\Latin{faciat}.  After impersonal negative \Latin{non est}, the clause characterizes the
nonexistent class of one who does good.

\textbf{Translation.}  ``There is no one who does good; there is not even one.''
\end{workedexample}

\begin{workedexample}{an actual identity (1 John 5:5)}{W.II.41}
\Latin{Quis est qui vincit mundum, nisi qui credit quoniam Iesus est Filius Dei?}

\textbf{Analysis.}  Both \Latin{vincit} and \Latin{credit} are indicative.  The question is
answered by an actual, asserted identity---the one who believes---rather than a merely
possible class.  \Latin{Qui} is nominative in each clause.

\textbf{Translation.}  ``Who is it that conquers the world, except the one who believes that
Jesus is the Son of God?''
\end{workedexample}

\begin{workedexample}{factual relative and jussive (Matthew 11:15)}{W.II.42}
\Latin{Qui habet aures audiendi, audiat.}

\textbf{Analysis.}  Indicative \Latin{habet} asserts the condition of an open class: whoever
in fact has ears for hearing.  Subjunctive \Latin{audiat} is an independent jussive, not part
of the relative clause.  Genitive gerund \Latin{audiendi} depends on \Latin{aures}.

\textbf{Translation.}  ``He who has ears for hearing, let him hear.''
\end{workedexample}

\begin{workedexample}{relative ground in the Gloria}{W.II.43}
\Latin{Domine Fili unigenite, Iesu Christe; Domine Deus, Agnus Dei, Filius Patris; qui tollis
peccata mundi, miserere nobis.}

\textbf{Analysis.}  \Latin{Qui} is nominative masculine singular and refers to the addressee
through the chain of vocative titles.  Indicative \Latin{tollis} states the known ground of
the plea; imperative \Latin{miserere} is the main request.

\textbf{Translation.}  ``Lord, only-begotten Son, Jesus Christ; Lord God, Lamb of God, Son of
the Father; you who take away the sins of the world, have mercy on us.''
\end{workedexample}

\section{Writing laboratory}

Copy \Latin{Non est qui faciat bonum} beside \Latin{Agnus qui tollit peccata mundi}.  Mark the
same morphology of \Latin{qui}, then explain in one sentence why the verb moods differ.

Compose a factual relative after \Latin{sacerdos}, a characteristic after \Latin{nemo est}, a
relative purpose after \Latin{misit servum}, and an identity question after \Latin{quis est}.
For the last, write one indicative and one subjunctive version and explain the different
viewpoint.



\end{courseunit}

\begin{courseunit}{II.11}
\chapter{Commands, prohibitions, wishes, and fearing}

\begin{learningobjectives}
  \item form present and future imperatives, including common irregulars;
  \item use jussive and hortatory subjunctives for first and third persons;
  \item distinguish the ordinary prohibition patterns;
  \item read fulfilled, possible, and unfulfilled wishes and the reversed negatives of fear.
\end{learningobjectives}

\section{Direct volition}

An \Term{imperative} directly issues a command, most commonly to a second
person. A \Term{jussive subjunctive} expresses a third-person command; a \Term{hortatory subjunctive}
urges first-person plural action (``let us''). A \Term{prohibition} is a
negative command. An \Term{optative subjunctive} expresses a wish, while a
\Term{fear clause} names the event whose occurrence or nonoccurrence is feared.
These labels identify distinct uses even when English renders several with
``may'' or ``let.''

\begin{grammarbox}{Imperatives and jussives}
\begin{center}
\begin{tabular}{@{}llllll@{}}
\toprule
& \Latin{amo} & \Latin{moneo} & \Latin{rego} & \Latin{capio} & \Latin{audio}\\ \midrule
2 sg. active & ama & mone & rege & cape & audi\\
2 pl. active & amate & monete & regite & capite & audite\\
2 sg. passive & amare & monere & regere & capere & audire\\
2 pl. passive & amamini & monemini & regimini & capimini & audimini\\
\bottomrule
\end{tabular}
\end{center}

Four frequent singular active forms drop \Latin{-e}: \Latin{dic, duc, fac, fer}; plurals are
\Latin{dicite, ducite, facite, ferte}.  Present subjunctive supplies first- and third-person
commands: \Latin{oremus}, ``let us pray'' (hortatory); \Latin{veniat}, ``let him come''
(jussive).  Future imperatives, common in law and rubrics, use \Latin{-to/-tote/-nto}; the
essential forms of \Latin{sum} are \Latin{esto, estote, sunto}.
\end{grammarbox}

\begin{grammarbox}{Prohibition, wish, and fear}
A direct second-person prohibition normally uses \Latin{noli/nolite} + infinitive:
\Latin{nolite timere}.  \Latin{Ne} + perfect subjunctive can give a sharp prohibition
(\Latin{ne timueris}); present subjunctive with \Latin{ne} is common in prayers and third
persons (\Latin{ne respicias}).

A possible or fulfilled wish uses present/perfect subjunctive, often with \Latin{utinam}; an
unfulfilled present wish uses imperfect, and an unfulfilled past wish pluperfect:
\Latin{utinam adsit}; \Latin{utinam adesset}; \Latin{utinam adfuisset}.  Liturgical
\Latin{fiat, sit, requiescat} is frequently optative or jussive without \Latin{utinam}.

After a verb of fearing, \Latin{ne} introduces what one fears will happen; \Latin{ut} or
\Latin{ne non} introduces what one fears will not happen.  Translate the fear, not the marker
word mechanically: \Latin{timeo ne veniat}, ``I fear that he may come'';
\Latin{timeo ut veniat}, ``I fear that he may not come.''
\end{grammarbox}

\section{Reading the forms in context}

\begin{workedexample}{institution commands in the Canon}{W.II.44}
\Latin{Accipite et manducate ex hoc omnes: hoc est enim Corpus meum.}

\textbf{Analysis.}  \Latin{Accipite} and \Latin{manducate} are second-person plural present
active imperatives.  \Latin{Omnes} is a substantival nominative vocative-like address; the
prepositional phrase \Latin{ex hoc} refers to what is received.

\textbf{Translation.}  ``All of you, take and eat of this, for this is my Body.''
\end{workedexample}

\begin{workedexample}{prayerful prohibition before the Peace}{W.II.45}
\Latin{Ne respicias peccata mea, sed fidem Ecclesiæ tuæ; eamque secundum voluntatem tuam
pacificare et coadunare digneris.}

\textbf{Analysis.}  Present subjunctives \Latin{respicias} and \Latin{digneris} express
negative and positive petitions.  \Latin{Ne} negates the first; \Latin{sed} supplies the
contrast.  Deponent \Latin{digneris} governs infinitives \Latin{pacificare et coadunare}.

\textbf{Translation.}  ``Look not upon my sins, but upon the faith of your Church; and be
pleased to give her peace and unity according to your will.''
\end{workedexample}

\begin{workedexample}{jussive prayer in the Pater noster}{W.II.46}
\Latin{Fiat voluntas tua, sicut in cælo et in terra.}

\textbf{Analysis.}  \Latin{Fiat} is third-person singular present subjunctive of \Latin{fio},
used independently as a jussive/optative.  Nominative \Latin{voluntas tua} is its subject.

\textbf{Translation.}  ``May your will be done, as in heaven and on earth.''
\end{workedexample}

\begin{workedexample}{unfulfilled wish (Apocalypse 3:15)}{W.II.47}
\Latin{Utinam frigidus esses aut calidus!}

\textbf{Analysis.}  \Latin{Utinam} marks a wish.  Imperfect subjunctive \Latin{esses} shows
that the desired present state is contrary to fact.

\textbf{Translation.}  ``Would that you were cold or hot!''
\end{workedexample}

\begin{workedexample}{fear of a completed possibility (Galatians 4:11)}{W.II.48}
\Latin{Timeo vos, ne forte sine causa laboraverim in vobis.}

\textbf{Analysis.}  \Latin{Ne} introduces what Paul fears may be true.  Perfect subjunctive
\Latin{laboraverim} is prior to present \Latin{timeo}; \Latin{forte} softens the feared
possibility.

\textbf{Translation.}  ``I fear for you, lest perhaps I have labored among you in vain.''
\end{workedexample}

\section{Writing laboratory}

Write the active and passive imperative grid, then add \Latin{dic, duc, fac, fer}.  Copy
\Latin{Accipite et manducate} and \Latin{Fiat voluntas tua}; label person and command type.

Give Latin for ``Let us pray,'' ``Let him receive,'' ``Do not fear'' (two patterns), ``May
they rest in peace,'' ``Would that he were here,'' and ``I fear that he will not come.''  For
the last item write both \Latin{ut} and \Latin{ne non} versions.



\end{courseunit}

\begin{courseunit}{II.12}
\chapter{Gerund, gerundive, supine, and passive periphrastic}

\begin{learningobjectives}
  \item decline the gerund and distinguish it from a gerundive agreeing with a noun;
  \item convert an object-taking gerund phrase to the preferred gerundive construction;
  \item recognize the two surviving supines;
  \item translate obligation and identify the dative of agent in a passive periphrastic.
\end{learningobjectives}

\section{Verbal nouns and adjectives}

\begin{grammarbox}{Gerund}
The gerund is a neuter singular verbal noun made from the present stem plus
\Latin{-nd-}: first and second conjugations show \Latin{-and-/-end-}, third \Latin{-end-},
and third-\Latin{-io}/fourth normally \Latin{-iend-}.  It has no nominative; the infinitive
serves there.

\begin{center}
\begin{tabular}{@{}lllll@{}}
\toprule
& Genitive & Dative (rare) & Accusative & Ablative\\ \midrule
\Latin{amo} & amandi & amando & ad amandum & amando\\
\Latin{rego} & regendi & regendo & ad regendum & regendo\\
\Latin{audio} & audiendi & audiendo & ad audiendum & audiendo\\
\bottomrule
\end{tabular}
\end{center}

The genitive depends on a noun/adjective (\Latin{tempus orandi}); accusative normally follows
\Latin{ad} for purpose; ablative expresses means or follows a preposition.  A gerund can govern
an object, but Latin normally avoids it by putting that object into the gerund's case and
making a gerundive agree: \Latin{ad offerendum hostiam} becomes the preferred
\Latin{ad hostiam offerendam}.
\end{grammarbox}

\begin{grammarbox}{Gerundive and passive periphrastic}
The gerundive is a first/second-declension verbal adjective:
\Latin{amandus, monendus, regendus, capiendus, audiendus}.  It agrees with its noun and has
passive force.  In an attributive purpose phrase, translate idiomatically:
\Latin{ad hostiam offerendam}, ``for offering the victim.''

Gerundive + a form of \Latin{sum} is the passive periphrastic of necessity:
\Latin{hostia offerenda est}, ``the victim must be offered.''  The person obliged is normally
dative: \Latin{sacerdoti hostia offerenda est}, ``the priest must offer the victim'' (more
literally, ``the victim is to-be-offered for the priest'').  With an intransitive verb the
construction is impersonal and neuter singular: \Latin{nobis orandum est}, ``we must pray.''
\end{grammarbox}

\begin{grammarbox}{Supine}
The supine is a fourth-declension verbal noun with only two live forms.  Accusative
\Latin{-um} may express purpose after motion (\Latin{venit oratum}, ``he came to pray'').
Ablative \Latin{-u} gives respect after certain adjectives (\Latin{mirabile dictu},
``wonderful to say'').  It takes an object as a verb can.  The construction is rare in the
Missal and Vulgate, which prefer an infinitive, \Latin{ad} + gerund(ive), or a purpose clause;
the reader's goal is recognition, not imitation everywhere.
\end{grammarbox}

\section{Reading the forms in context}

\begin{workedexample}{genitive gerund (Ecclesiastes 3:2)}{W.II.49}
\Latin{Tempus nascendi, et tempus moriendi.}

\textbf{Analysis.}  \Latin{Nascendi} and \Latin{moriendi} are genitive gerunds depending on
\Latin{tempus}.  Their deponent dictionary forms do not make these gerunds passive.

\textbf{Translation.}  ``A time for being born and a time for dying.''
\end{workedexample}

\begin{workedexample}{purpose gerund with an object (Psalm 91:3)}{W.II.50}
\Latin{Ad annuntiandum mane misericordiam tuam, et veritatem tuam per noctem.}

\textbf{Analysis.}  Accusative gerund \Latin{annuntiandum} follows \Latin{ad} and governs
objects \Latin{misericordiam} and \Latin{veritatem}.  Regularized gerundive Latin would be
\Latin{ad misericordiam tuam et veritatem tuam annuntiandas}.

\textbf{Translation.}  ``To proclaim your mercy in the morning and your truth through the
night.''
\end{workedexample}

\begin{workedexample}{ablative gerund (Romans 15:13)}{W.II.51}
\Latin{Deus autem spei repleat vos omni gaudio et pace in credendo.}

\textbf{Analysis.}  \Latin{Credendo} is ablative gerund after \Latin{in}, ``in believing.''
Present subjunctive \Latin{repleat} is an independent prayer; \Latin{vos} is its object.

\textbf{Translation.}  ``May the God of hope fill you with all joy and peace in believing.''
\end{workedexample}

\begin{workedexample}{necessity without a gerundive (1 Corinthians 9:16)}{W.II.52}
\Latin{Necessitas enim mihi incumbit; væ enim mihi est, si non evangelizavero.}

\textbf{Analysis.}  Dative \Latin{mihi} marks the person under necessity, the same role a
dative of agent has in \Latin{mihi evangelizandum est}.  This authentic phrasing therefore
provides a safe transformation model without pretending the Vulgate used a periphrastic here.

\textbf{Translation.}  ``For necessity lies upon me; indeed, woe is me if I do not preach the
Gospel.''
\end{workedexample}

\begin{workedexample}{the Missal's preferred purpose infinitive in the Creed}{W.II.53}
\Latin{Et iterum venturus est cum gloria, iudicare vivos et mortuos.}

\textbf{Analysis.}  \Latin{Venturus est} is the future active periphrastic, ``is going to
come,'' not a gerundive of obligation.  Present infinitive \Latin{iudicare} expresses purpose
where earlier Latin might use a supine after motion.

\textbf{Translation.}  ``And he will come again with glory to judge the living and the dead.''
\end{workedexample}

\section{Writing laboratory}

Write the gerund paradigm for all five conjugations.  Beside every form write a gerundive
phrase in the same case: \Latin{orandi / orationis faciendæ}; \Latin{ad offerendum / ad
hostiam offerendam}.  Circle the noun that forces gerundive agreement.

Convert \Latin{ad laudandum Deum} to \Latin{ad Deum laudandum}; write ``The gifts must be
blessed by the priest'' with dative agent; write ``We must pray'' impersonally; change
\Latin{venit oratum} into an infinitive and then a purpose clause.  Explain why
\Latin{venturus} is not \Latin{veniendus}.



\end{courseunit}

\begin{courseunit}{II.13}
\chapter{Place, time, means, manner, and cause}

\begin{learningobjectives}
  \item distinguish place where, motion to, motion from, and route by which;
  \item distinguish time when, time within which, and duration;
  \item separate instrument, personal agent, accompaniment, and manner;
  \item identify source, cause, and motive from both case and meaning.
\end{learningobjectives}

\section{The adverbial cases}

An \Term{adverbial case use} is a noun phrase that modifies the circumstances
of an action---where, when, how, why, or by what means---rather than serving as
subject or direct object. The case narrows the possibilities, but the verb,
preposition, and meaning determine the exact relation.

\begin{grammarbox}{Place and time}
\begin{tabularx}{\linewidth}{@{}P{1.35in}P{1.5in}Y@{}}
\toprule
\textbf{Relation} & \textbf{Normal construction} & \textbf{Example}\\ \midrule
place where & \Latin{in} + ablative & \Latin{in terra}; towns, small islands,
\Latin{domus}, and \Latin{rus} may use locative: \Latin{Romæ, domi, ruri}\\
motion to & \Latin{in/ad} + accusative & \Latin{ad altare, in cælum}; towns and
\Latin{domum/rus} normally omit preposition: \Latin{Romam, domum}\\
motion from & \Latin{a/ab, e/ex, de} + ablative & \Latin{de cælis, ex altari}; towns and
\Latin{domus/rus} may use bare ablative: \Latin{Roma, domo}\\
route by which & ablative, often \Latin{per} + accusative & \Latin{eadem via};
\Latin{per viam} emphasizes passage through\\
time when/within & ablative without preposition & \Latin{tertia die}; \Latin{tribus diebus},
``within three days,'' when context requires\\
duration/extent & accusative, often \Latin{per} & \Latin{tres dies};
\Latin{per omnia sæcula}\\
\bottomrule
\end{tabularx}
\end{grammarbox}

\begin{grammarbox}{Means, agent, manner, and cause}
\Term{Means} or \Term{instrument} is the nonpersonal thing by which an action
is accomplished and normally uses the bare ablative: \Latin{verbo sanat}, ``he heals by a
word.''  A personal agent with a passive verb uses \Latin{a/ab}: \Latin{a sacerdote
offertur}. \Term{Accompaniment} identifies who or what is together with the
subject; \Term{manner} describes how the action is performed. Accompaniment
uses \Latin{cum}; manner uses \Latin{cum} + ablative, but may omit
\Latin{cum} when an adjective makes the manner clear: \Latin{magna voce}, \Latin{pura mente}.

Bare ablative may express cause (\Latin{timore tremit}) or source; \Latin{propter/ob} +
accusative normally gives motive, \Latin{causa/gratia} + preceding genitive gives purpose,
and \Latin{de/ex} + ablative marks source.  Do not infer a label from the preposition alone:
the governing verb and sense decide.
\end{grammarbox}

\section{Reading the forms in context}

\begin{workedexample}{direction and location at the opening prayers}{W.II.54}
\Latin{Introibo ad altare Dei, ad Deum qui lætificat iuventutem meam.}

\textbf{Analysis.}  Both instances of \Latin{ad} govern accusative and mark direction toward,
not location at.  The second phrase appositionally interprets the altar's goal as God.

\textbf{Translation.}  ``I will go in to the altar of God, to God who gives joy to my youth.''
\end{workedexample}

\begin{workedexample}{two locations in the Gloria}{W.II.55}
\Latin{Gloria in excelsis Deo. Et in terra pax hominibus bonæ voluntatis.}

\textbf{Analysis.}  \Latin{In excelsis} and \Latin{in terra} are place where with ablative.
\Latin{Deo} and \Latin{hominibus} are datives of reference or recipient; forms must not be
relabeled ablative merely because a nearby phrase is ablative.

\textbf{Translation.}  ``Glory to God in the highest places, and on earth peace to people of
good will.''
\end{workedexample}

\begin{workedexample}{source and motive in the Creed}{W.II.56}
\Latin{Qui propter nos homines et propter nostram salutem descendit de cælis.}

\textbf{Analysis.}  Repeated \Latin{propter} + accusative marks beneficiary/motive;
\Latin{de cælis} marks motion from a higher source.  Nominative \Latin{qui} is subject of
\Latin{descendit}.

\textbf{Translation.}  ``Who for us human beings and for our salvation came down from
heaven.''
\end{workedexample}

\begin{workedexample}{time when and motion into in the Creed}{W.II.57}
\Latin{Et resurrexit tertia die secundum Scripturas, et ascendit in cælum.}

\textbf{Analysis.}  \Latin{Tertia die} is bare ablative of time when.  \Latin{In cælum} is
accusative of motion into; it does not mean that the ascension happened while already in
heaven.

\textbf{Translation.}  ``And he rose on the third day according to the Scriptures, and
ascended into heaven.''
\end{workedexample}

\begin{workedexample}{means and manner in a Postcommunion}{W.II.58}
\Latin{Quod ore sumpsimus, Domine, pura mente capiamus.}

\textbf{Analysis.}  Bare ablative \Latin{ore} is bodily means; adjective + ablative
\Latin{pura mente} expresses manner without \Latin{cum}.  The parallel English preposition
``with'' must not hide their different functions.

\textbf{Translation.}  ``Lord, may we receive with a pure mind what we have taken by mouth.''
\end{workedexample}

\section{Writing laboratory}

Draw four arrows labelled \textsc{where / to / from / by way of} and place
\Latin{in ecclesia, ad ecclesiam, ex ecclesia, eadem via} beneath them.  Add the exceptional
series \Latin{Romæ, Romam, Roma} and \Latin{domi, domum, domo}.

Compose paired contrasts: ``at the altar / to the altar''; ``on the third day / for three
days''; ``by the priest / by a word''; ``with the disciples / with great faith''; ``from
heaven / for our salvation.''  Parse every noun before translating.



\end{courseunit}

\begin{courseunit}{II.14}
\chapter{Word order, ellipsis, attraction, and liturgical compression}

\begin{learningobjectives}
  \item reconstruct clause structure when related words are separated;
  \item supply only ellipses licensed by grammar and parallelism;
  \item recognize incorporated antecedents and mood/case attraction;
  \item unpack the address--ground--petition architecture of compressed liturgical prose.
\end{learningobjectives}

\section{Reading what endings connect}

\Term{Word order} is the sequence in which words appear; \Term{syntax} is the
network of grammatical relations among them. Latin can separate words that
belong together, so changing the order need not change the syntax. Always use
endings, agreement, and government to recover that network before relying on
position.

\begin{grammarbox}{Order and emphasis}
Latin commonly places subject--object--verb, modifier near noun, and governing word after its
dependents, but none is an inviolable rule.  Initial and final positions carry emphasis;
pronouns and particles often come early; verbs can be delayed to close a period.  Hyperbaton
separates agreeing words around material they frame: \Latin{sanctas \ldots{} manus suas}.
Chiasmus arranges ABBA; asyndeton omits conjunctions; anaphora repeats a word.  Recover syntax
from endings and government before rearranging into English.
\end{grammarbox}

\begin{grammarbox}{Licensed ellipsis}
\Term{Ellipsis} is the omission of a word that grammar and context make
recoverable. Latin readily omits a form of \Latin{sum}, a repeated verb, an antecedent recoverable from a
relative, or a noun recoverable from an adjective/participle.  Supply the smallest element that
makes the attested syntax complete: \Latin{Sanctus Dominus} can imply \Latin{est};
\Latin{viventes Deo} can share prior \Latin{esse}.  Do not insert an agent, causal relation, or
theological interpretation that the grammar leaves open.
\end{grammarbox}

\begin{grammarbox}{Attraction and incorporation}
\Term{Attraction} is the adjustment of a form or mood under the influence of a
nearby construction rather than its simplest independent syntax. A relative's
normal case comes from its own clause, but Latin sometimes attracts it into the
case of its antecedent, especially when the antecedent itself is omitted or incorporated:
\Latin{in qua mensura}, ``in the measure in which.''  A relative or subordinate clause inside
subjunctive discourse may itself take subjunctive by \Term{modal attraction}, especially when
it belongs to the thought rather than the narrator's assertion.  Attraction is a diagnosis
after ordinary agreement and government fail, not a shortcut around parsing.
\end{grammarbox}

\begin{grammarbox}{The compressed collect and Canon period}
Use this recovery order:
\begin{enumerate}
  \item circle every finite verb and conjunction;
  \item locate vocative address and nominative subject;
  \item attach each accusative/dative to a governing verb;
  \item match separated adjectives, participles, and relatives by endings;
  \item bracket subordinate clauses from marker to finite verb;
  \item supply parallel ellipses, then make a literal prose-order version;
  \item only then produce readable English.
\end{enumerate}
Liturgical prose often has one finite petition supporting many vocatives, appositions,
participles, genitives, and prepositional phrases.  Punctuation is a guide, not a substitute
for these steps.
\end{grammarbox}

\section{Reading the forms in context}

\begin{workedexample}{fronting and delayed government in the Canon}{W.II.59}
\Latin{Te igitur, clementissime Pater, per Iesum Christum Filium tuum Dominum nostrum,
supplices rogamus ac petimus.}

\textbf{Analysis.}  Accusative \Latin{te} is fronted far before its governing
\Latin{rogamus ac petimus}.  Vocative \Latin{Pater} is expanded by \Latin{clementissime};
\Latin{Filium} and \Latin{Dominum} are appositions to \Latin{Iesum Christum}.  Nominative
\Latin{supplices} describes the implicit subject \Latin{nos}.

\textbf{Translation.}  ``Therefore, most merciful Father, through Jesus Christ your Son our
Lord, we humbly ask and petition you.''
\end{workedexample}

\begin{workedexample}{an adjective chain and final verb in the Canon}{W.II.60}
\Latin{Quam oblationem tu, Deus, in omnibus, quæsumus, benedictam, adscriptam, ratam,
rationabilem, acceptabilemque facere digneris.}

\textbf{Analysis.}  Accusative \Latin{quam oblationem} is the object of \Latin{facere}; five
accusative feminine adjectives are its object complements.  Deponent \Latin{digneris} governs
\Latin{facere} and comes last; parenthetic \Latin{quæsumus} carries the plea.  Enclitic
\Latin{-que} joins the final adjective to the series.

\textbf{Translation.} ``Be pleased, O God, we ask, to make this offering
blessed in every respect, enrolled, ratified, rational/reasonable, and
acceptable.'' The conventional liturgical rendering \English{spiritual} is
theological and idiomatic rather than the closest lexical gloss.
\end{workedexample}

\begin{workedexample}{ellipsis and acclamation in the Sanctus}{W.II.61}
\Latin{Sanctus, Sanctus, Sanctus Dominus Deus Sabaoth. Pleni sunt cæli et terra gloria tua.}

\textbf{Analysis.} The triple nominative adjectives form a verbless
acclamation. Analytic English may supply \Latin{[est]}, but Latin does not
require a predicate-nominative parse: the adjectives are predicative or
acclamatory modifiers of \Latin{Dominus Deus}. In the next sentence,
\Latin{pleni} agrees with \Latin{cæli et terra}, and ablative
\Latin{gloria tua} expresses fullness.

\textbf{English rendering status: unavailable.} Continue from the Latin and
grammatical analysis without a model English text.
\end{workedexample}

\begin{workedexample}{incorporated antecedent (Matthew 7:2)}{W.II.62}
\Latin{In quo enim iudicio iudicaveritis, iudicabimini; et in qua mensura mensi fueritis,
remetietur vobis.}

\textbf{Analysis.}  \Latin{Quo iudicio} and \Latin{qua mensura} incorporate the standards
into their relative clauses: ``with the judgment with which,'' ``in the measure in which.''
Future perfect \Latin{iudicaveritis} and \Latin{mensi fueritis} precede the
main futures. \Latin{iudicabimini} is second-person plural future passive of
\Latin{iudico}; \Latin{remetietur} is formally third-person singular future
indicative of deponent \Latin{remetior, remetiri, remensus sum}, exceptionally
used with passive sense here, \English{will be measured back}.

\textbf{Translation.}  ``For with the judgment with which you judge, you will be judged; and
in the measure in which you have measured, it will be measured back to you.''
\end{workedexample}

\begin{workedexample}{nested compression in the Canon}{W.II.63}
\Latin{Hanc igitur oblationem servitutis nostræ, sed et cunctæ familiæ tuæ, quæsumus,
Domine, ut placatus accipias.}

\textbf{Analysis.}  Fronted \Latin{hanc oblationem} is the object of \Latin{accipias}; its two
genitives identify those whose offering it is.  Parenthetic \Latin{quæsumus} introduces the
substantive \Latin{ut}-request.  Predicate participle
\Latin{placatus} agrees with the implicit subject \Latin{tu}: ``being appeased/gracious.''

\textbf{Translation.}  ``Therefore we ask you, Lord, graciously to accept this offering of our
service and also of your whole household.''
\end{workedexample}

\section{Writing laboratory}

Copy the \Latin{Quam oblationem} sentence on alternate lines.  Connect every adjective to
\Latin{oblationem}, draw an arrow from \Latin{facere} to the object, and another from
\Latin{digneris} to \Latin{facere}.  Beneath it write a prose-order Latin version, retaining
every word.

Do the same with a collect chosen at random: (1) diplomatic copy, (2) marked syntax,
(3) prose-order Latin, (4) close English, (5) idiomatic English.  Then compress the composed
sentence \Latin{Rogamus te, Pater clementissime, ut dona quæ offerimus benedicas} by fronting
\Latin{te}, inserting \Latin{supplices}, and delaying the two governing verbs.

\end{courseunit}
